% ==========================================================
% แผนบริหารการสอนรายวิชาควอนตัมฟิสิกส์เบื้องต้น (Course Syllabus)
% ตามกรอบมาตรฐานผลลัพธ์การเรียนรู้ (OBE) มหาวิทยาลัยราชภัฏรำไพพรรณี
% ==========================================================

\chapter*{แผนบริหารการสอนรายวิชา}
\addcontentsline{toc}{chapter}{แผนบริหารการสอนรายวิชา}

\section*{ข้อมูลทั่วไปของรายวิชา}
\noindent\textbf{รหัสและชื่อรายวิชา} ฟิสิกส์ควอนตัมเบื้องต้น (Introduction to Quantum Physics)\\
\noindent\textbf{จำนวนหน่วยกิต} 3(3-0-6) บรรยาย 3 ชั่วโมง ศึกษาค้นคว้าด้วยตนเอง 6 ชั่วโมงต่อสัปดาห์\\
\noindent\textbf{หลักสูตรและประเภทของรายวิชา} วิทยาศาสตรบัณฑิต หมวดวิชาเฉพาะด้านฟิสิกส์ประยุกต์และดิจิทัลการเกษตร\\
\noindent\textbf{อาจารย์ผู้รับผิดชอบรายวิชาและผู้สอน} ผู้ช่วยศาสตราจารย์ ดร.ชีวะ ทัศนา\\
\noindent\textbf{ภาคการศึกษาและชั้นปีที่เรียน} ภาคการศึกษาที่ 1 ชั้นปีที่ 3

\vspace{10pt}
\section*{คำอธิบายรายวิชา}
การกำเนิดของกลศาสตร์ควอนตัม ปรากฏการณ์การแผ่รังสีของวัตถุดำ โฟโตอิเล็กทริก และการกระเจิงคอมป์ตัน ความเป็นทวิภาพของคลื่นและอนุภาค คลื่นสสารเดอ บรอยล์ กลุ่มคลื่นเกาส์เซียน และหลักความไม่แน่นอนของไฮเซนเบิร์ก ปริภูมิฮิลเบิร์ต สัญกรณ์บรา-เคทของดิแรก ตัวดำเนินการแฮร์มิเชียน สมการชเรอดิงเงอร์ที่ขึ้นและไม่ขึ้นกับเวลา การประยุกต์ในบ่อศักย์ กำแพงศักย์ และปรากฏการณ์ทันเนลลิง ตัวแกว่งกวัดฮาร์มอนิกควอนตัมและพีชคณิตตัวดำเนินการบันได อะตอมไฮโดรเจนในระบบพิกัดทรงกลม โมเมนตัมเชิงมุม สปินของอิเล็กตรอนและเมทริกซ์เพาลี ทฤษฎีบทการรบกวน สถิติควอนตัม สารสนเทศควอนตัม คิวบิต ทรงกลมบล็อค เกตควอนตัม ความพัวพันเชิงควอนตัม โพรโทคอลการส่งผ่านสถานะควอนตัมทางไกล อัลกอริทึมควอนตัมเบื้องต้น การจำลองเสมือนจริงสามมิติผ่านเทคโนโลยีภาพเสมือนผสานโลกจริง (AR/XR) และการประยุกต์โครงข่ายประสาทเทียมตามหลักฟิสิกส์ (Physics-Informed Neural Networks) ในการคำนวณเชิงตัวเลข

\vspace{10pt}
\section*{ผลลัพธ์การเรียนรู้ระดับรายวิชา (Course Learning Outcomes: CLOs)}
เมื่อสิ้นสุดการเรียนการสอนรายวิชานี้ นักศึกษาสามารถ
\begin{enumerate}[leftmargin=1.5cm]
    \item \textbf{CLO 1} อธิบายและวิเคราะห์หลักการพื้นฐานทางกลศาสตร์ควอนตัม สัญกรณ์ดิแรก และการแปลงสภาพเชิงเอกภาพได้อย่างถูกต้องตามหลักการทางคณิตศาสตร์และฟิสิกส์
    \item \textbf{CLO 2} แก้สมการชเรอดิงเงอร์สำหรับระบบศักย์หนึ่งมิติ ตัวแกว่งกวัดฮาร์มอนิก และอะตอมไฮโดรเจน พร้อมทั้งตีความฟังก์ชันคลื่นและความน่าจะเป็นได้อย่างถูกต้อง
    \item \textbf{CLO 3} วิเคราะห์และออกแบบวงจรสารสนเทศควอนตัม ระบบคิวบิตบนทรงกลมบล็อค สถานะพัวพันเบลล์ และโพรโทคอลโทรเลขควอนตัมทางไกลได้อย่างเป็นระบบ
    \item \textbf{CLO 4} พัฒนาโปรแกรมภาษา Python ประยุกต์ใช้ Qiskit, QuTiP และโครงข่ายประสาทเทียมตามหลักฟิสิกส์ (PINNs) เพื่อจำลองและหาผลเฉลยของระบบควอนตัมได้อย่างมีประสิทธิภาพ
    \item \textbf{CLO 5} ใช้งานสื่อจำลองเสมือนจริงสามมิติ WebAR ผ่านเทคโนโลยีการติดตามภาพ (Image Tracking) และรหัสตอบสนองเร็ว (QR Code) เพื่อการทดลองเชิงประจักษ์แบบไร้สัมผัส
\end{enumerate}

\newpage
\section*{แผนผังโครงสร้างเนื้อหาและนวัตกรรมการเรียนรู้ประจำบท (Master Syllabus Matrix)}

\begin{table}[H]
\centering
\small
\begin{tabularx}{\textwidth}{c l Y c}
\toprule
\textbf{บทที่} & \textbf{ชื่อบทเรียนภาษาไทย} & \textbf{นวัตกรรมดิจิทัล: AR Target, QR Lab \& Python} & \textbf{CLO} \\
\midrule
1 & สมบัติเชิงอนุภาคของรังสีและโฟตอนิกส์ยุคใหม่ & QR Lab 1: สเปกตรัมวัตถุดำและโฟตอนเดี่ยว\newline Python: Curve Fitting กฎพลังค์เทียบวีน & 1, 4 \\
\addlinespace
2 & ความเป็นทวิภาคคลื่น-อนุภาคและการแทรกสอด & AR Target: ภาพเครื่องมือสลิตคู่อิเล็กตรอน\newline Python: การจำลองการบานออกของ Wave Packet & 1, 4, 5 \\
\addlinespace
3 & สมการชเรอดิงเงอร์และพีชคณิตเชิงเส้นควอนตัม & QR Lab 3: ตัวดำเนินการและการแปลงฐาน Dirac\newline Python: เมทริกซ์แฮร์มิเชียนและค่าไอเกน & 1, 2, 4 \\
\addlinespace
4 & กลศาสตร์คลื่นในบ่อศักย์และปรากฏการณ์ทันเนลลิง & AR Target: ภาพหัวเข็ม STM ระดับอะตอม\newline Python PINNs: แก้บ่อศักย์และจำลองทันเนลลิง & 2, 4, 5 \\
\addlinespace
5 & ตัวแกว่งกวัดฮาร์มอนิกควอนตัมและพีชคณิตตัวดำเนินการ & QR Lab 5: ฟังก์ชัน Wigner และ Coherent State\newline Python QuTiP: ตัวดำเนินการบันได $a, a^\dagger$ & 2, 4 \\
\addlinespace
6 & ระบบอะตอมในสามมิติ โมเมนตัมเชิงมุม และสปิน & AR Target: ภาพออร์บิทัล $3d$ และ $4f$ สามมิติ\newline 3D WebGL: Orbital Viewer หมุนอิสระรอบแกน & 2, 5 \\
\addlinespace
7 & ทฤษฎีการรบกวน สถิติควอนตัม และสสารควอนตัม & QR Lab 7: การแยกของระดับพลังงาน Stark Effect\newline Python: สถิติ Fermi-Dirac และ Bose-Einstein & 2, 4 \\
\addlinespace
8 & สารสนเทศควอนตัม คิวบิต และทรงกลมบล็อค & AR Target: แผนภาพทรงกลมบล็อคสามมิติ\newline Python Qiskit: วงจร 1 คิวบิตและ Statevector & 3, 4, 5 \\
\addlinespace
9 & ความพัวพันเชิงควอนตัมและการสื่อสารควอนตัม & QR Lab 9: การวัดสถานะเบลล์และโทรเลขควอนตัม\newline QR Lab 12: การแจกจ่ายกุญแจรหัสลับควอนตัม QKD\newline Python Qiskit: Teleportation \& QKD BB84 & 3, 4, 5 \\
\addlinespace
10 & อัลกอริทึมควอนตัม คอมพิวติง และพรมแดนประยุกต์ & QR Lab 10: สถาปัตยกรรมชิปควอนตัมและโกรเวอร์\newline QR Lab 11: การแปลงฟูเรียร์ควอนตัม QFT และชอร์\newline Python: Grover Search \& Shor's Factoring & 3, 4, 5 \\
\bottomrule
\end{tabularx}
\caption{โครงสร้างเนื้อหา ผลลัพธ์การเรียนรู้ และนวัตกรรมดิจิทัลประจำบทเรียน}
\end{table}

\vspace{8pt}
\section*{แผนการประเมินผลการเรียนรู้}
การวัดและประเมินผลการเรียนรู้ใช้ระบบการประเมินอิงเกณฑ์ตามมาตรฐานมหาวิทยาลัยราชภัฏรำไพพรรณี แบ่งสัดส่วนคะแนนดังนี้
\begin{enumerate}[leftmargin=1.5cm]
    \item \textbf{การมีส่วนร่วมและการทดลองในชั้นเรียนเสมือนจริง (Interactive Lab Participation)} ร้อยละ 15
    \item \textbf{รายงานปฏิบัติการคอมพิวเตอร์และการเขียนโปรแกรม Python (Computational Lab Reports)} ร้อยละ 20
    \item \textbf{แบบฝึกหัดและการแก้โจทย์ปัญหาเชิงวิเคราะห์แบบ P-S-C-T ประจำบท} ร้อยละ 15
    \item \textbf{การสอบประเมินผลกลางภาคการศึกษา (Midterm Examination)} ร้อยละ 25
    \item \textbf{การสอบประเมินผลปลายภาคการศึกษา (Final Examination)} ร้อยละ 25
\end{enumerate}
